\documentclass[11pt]{article}

\usepackage[utf8]{inputenc}
\usepackage[T1]{fontenc}
\usepackage{amsmath, amssymb, amsthm, mathtools}
\usepackage{braket}
\usepackage{geometry}
\usepackage{hyperref}
\usepackage{enumitem}

\geometry{margin=1in}

\newcommand{\outact}[2]{#1!\!\left(#2\right)}
\newcommand{\inact}[3]{\mathbf{for}\!\left(#1 \leftarrow #2\right)\!\left\{#3\right\}}
\newcommand{\Proc}{\mathit{Proc}}
\newcommand{\CTRL}{\mathrm{CTRL}}
\newcommand{\CNOT}{\mathrm{CNOT}}

\title{The MQ-Calculus:\\[4pt]
\large A Calculus for Quantum Processes Involving Intermittent Measurement}
\author{Mike Stay \and Lucius Gregory Meredith}
\date{February 23, 2026}

\begin{document}
\maketitle

%% ----------------------------------------------------------------
\section{Introduction and Motivation}
%% ----------------------------------------------------------------

We introduce a calculus for processes that involve measurements
interspersed throughout the computation.  The key intuition is that
\emph{communication = measurement}.

%% ----------------------------------------------------------------
\section{Related Work}
%% ----------------------------------------------------------------

The calculus is related to the CQP calculus of Gay and Nagarajan.  The
primary difference is that their calculus has an explicit measurement
operator, whereas this calculus insists that measurement happens when a
quantum state is communicated from one agent to another.

%% ----------------------------------------------------------------
\section{Technical Presentation}
%% ----------------------------------------------------------------

\subsection{Syntax}

\[
  \theta \;::=\; [0, 2\pi)
\]

\paragraph{States.}
\begin{align*}
  S, T \;::=\; & \ket{0} \;\mid\; \ket{1} \;\mid\; v \\
    \mid\; & H(S) \;\mid\; F(\theta)(S) \;\mid\; R(\theta)(S) \\
    \mid\; & \CTRL\text{-}F(\theta)(S, T)
            \;\mid\; \CNOT(S, T)
            \;\mid\; \CTRL(U)(S, S, S, \ldots)(T, T, T, \ldots) \\
    \mid\; & S \otimes T
            \;\mid\; \mathbb{C}\, S
            \;\mid\; S + T
\end{align*}

\paragraph{Processes.}
\begin{align*}
  Q &\;::=\; \bigotimes_{j}\, x_j!\!\left(S\right)
  &&\text{(output; when $j$ ranges over an empty domain this is $\mathbf{0}$)}\\[6pt]
  P &\;::=\; \mathbf{for}\!\left(\bigotimes_{j} m_j
             \leftarrow \bigotimes_{j} x_j \right)\!\left\{\Proc\right\}
  &&\text{(input; when $j$ ranges over an empty domain this is $\mathbf{0}$)}\\[6pt]
  \Proc &\;::=\; Q \;\mid\; P
         \;\mid\; \textstyle\sum_{i}\, p_i\;\Proc
         \;\mid\; \Proc \mid \Proc
  &&\text{where } p_i \in \mathbb{R}^{+}
\end{align*}

\subsection{Hypothesis}

\[
  p\;(P_1 \mid P_2) \;=\; p\,P_1 \;\mid\; p\,P_2
\]

\subsection{Distributing Sums over Parallel Composition}

Note that we have
\[
  \sum_{i_1} p_{1,i_1}
    \!\left(\bigotimes_{j} x_{i_1,j}!\!\left(S_{i_1}\right)\right)
  \;\Big|\;
  \sum_{i_2} p_{2,i_2}
    \!\left(
      \mathbf{for}\!\left(
        \bigotimes_{j} m_{i_2,j}
        \leftarrow
        \bigotimes_{j} x_{i_2,j}
      \right)\!\left\{\Proc_{i_2}\right\}
    \right)
\]
\[
  =\;
  \sum_{i_1, i_2} p_{1,i_1}\, p_{2,i_2}\;
    \left(\bigotimes_{j} x_{i_1,j}!\!\left(S_{i_1}\right)\right)
  \;\Big|\;
    \left(
      \mathbf{for}\!\left(
        \bigotimes_{j} m_{i_2,j}
        \leftarrow
        \bigotimes_{j} x_{i_2,j}
      \right)\!\left\{\Proc_{i_2}\right\}
    \right)
\]
This will be useful in writing the full comm rule.

\subsection{The Comm Rule}

We build the comm rule in steps.  The driving intuition is that
\[
  p_1 \;\bigotimes_{j \in J} x_j!\!\left(S\right)
  \;\Big|\;
  p_2 \;\mathbf{for}\!\left(
    \bigotimes_{j \in J} m_j \leftarrow \bigotimes_{j \in J} x_j
  \right)\!\left\{\Proc\right\}
  \;\;\longrightarrow\;\;
  p_1\, p_2 \;\sum_{r} \left|\braket{r | S}\right|^{2}\;
  \Proc\!\left\{r / m\right\}
\]
where $S$ is a $|J|$-qubit state.

\bigskip

Thus the comm rule in full is:

\bigskip

\noindent\textbf{Comm:}
\[
  \frac{\displaystyle
    \sum_{i_1} p_{1,i_1}
      \!\left(\bigotimes_{j_1} x_{i_1,j_1}!\!\left(S_{i_1}\right)\right)
    \;\Big|\;
    \sum_{i_2} p_{2,i_2}
      \!\left(
        \mathbf{for}\!\left(
          \bigotimes_{j_2} m_{i_2,j_2}
          \leftarrow
          \bigotimes_{j_2} x_{i_2,j_2}
        \right)\!\left\{\Proc_{i_2}\right\}
      \right)
  }{\displaystyle
    \sum_{\substack{i_1, i_2 \\[2pt]
      \bigotimes_j x_{i_1,j} \,=\, \bigotimes_j x_{i_2,j}}}
    p_{1,i_1}\, p_{2,i_2}
    \sum_{r}\;
    \left|\!\braket{r \,|\, S_{i_1}}\!\right|^{2}\;
    \Proc_{i_2}\!\left\{r \,/\, m_{i_2}\right\}
  }
\]

\noindent
where the outer sum ranges over all pairs $(i_1, i_2)$ such that the
channel tuples match, i.e.\
$\bigotimes_j x_{i_1, j} = \bigotimes_j x_{i_2, j}$, and the inner
sum over $r$ ranges over the $|J|$-qubit computational basis states.
For each such pair, the probability weight is the product
$p_{1,i_1}\, p_{2,i_2}$, the state $S_{i_1}$ is measured in the
computational basis yielding outcome $r$ with probability
$|\!\braket{r | S_{i_1}}\!|^2$, and the continuation
$\Proc_{i_2}\{r/m_{i_2}\}$ is the body of the input with the
measurement outcome substituted for the bound pattern variables.

%% ----------------------------------------------------------------
\section{Examples}
%% ----------------------------------------------------------------

\begin{enumerate}[label=\arabic*.]
  \item Simple communication:
  \[
    x!\!\left(S\right)
    \;\Big|\;
    \mathbf{for}\!\left(s \leftarrow x\right)\!\left\{P\right\}
  \]

  \item Iterated Hadamard:
  \[
    x!\!\left(\ket{0}\right)
    \;\Big|\;
    \mathbf{for}\!\left(s \leftarrow x\right)\!\left\{
      x!\!\left(Hs\right)
      \;\Big|\;
      \mathbf{for}\!\left(s \leftarrow x\right)\!\left\{
        x!\!\left(Hs\right)
      \right\}
    \right\}
  \]

  \item Collapsed form:
  \[
    x!\!\left(\ket{0}\right)
    \;\Big|\;
    \mathbf{for}\!\left(s \leftarrow x\right)\!\left\{
      x!\!\left(HHs\right)
    \right\}
  \]

  \item General circuit form:
  \[
    \mathit{wires}!\!\left(\ket{0\cdots 0}\right)
    \;\Big|\;
    \mathbf{for}\!\left(\mathit{initialState}
      \leftarrow \mathit{wires}\right)\!\left\{
      \mathit{output}!\!\left(U_{\mathrm{circuit}}\;\mathit{initialState}\right)
    \right\}
  \]
\end{enumerate}

\subsection{Notational Conventions}

\begin{align*}
  1\; x!\!\left(0\right) \otimes 1\; y!\!\left(0\right)
    &= 1\; (xy)!\!\left(00\right) \\[6pt]
  (H \otimes I)\bigl(1\; x!\!\left(0\right) \otimes 1\; y!\!\left(0\right)\bigr)
    &= \tfrac{1}{\sqrt{2}}\; (xy)!\!\left(00\right)
     + \tfrac{1}{\sqrt{2}}\; (xy)!\!\left(10\right) \\[6pt]
  \CNOT\!\left(\tfrac{1}{\sqrt{2}}\; (xy)!\!\left(00\right)
     + \tfrac{1}{\sqrt{2}}\; (xy)!\!\left(10\right)\right)
    &= \tfrac{1}{\sqrt{2}}\; (xy)!\!\left(00\right)
     + \tfrac{1}{\sqrt{2}}\; (xy)!\!\left(11\right) \\[6pt]
  \CNOT\!\left((H \otimes I)\bigl(1\; x!\!\left(0\right)
     &\otimes 1\; y!\!\left(0\right)\bigr)\right)
    \;\Big|\;
    \mathbf{for}\!\left(v \leftarrow x\right)\!\left\{X\right\}
    \;\Big|\;
    \mathbf{for}\!\left(w \leftarrow y\right)\!\left\{Y\right\}
\end{align*}

%% ----------------------------------------------------------------
\section{Analysis of Wigner's Friend Paradox}
%% ----------------------------------------------------------------

TBD

%% ----------------------------------------------------------------
\section{Alternative Design Options}
%% ----------------------------------------------------------------

\subsection{Classical Processes Acting on Quantum States}

A configuration consists of:
\begin{enumerate}
  \item A number $n$ of allocated wires.
  \item A $2^n$-dimensional state $\ket{\psi}$.
\end{enumerate}

A \emph{quantum process} acts on a configuration.

\subsection{De Bruijn Indices}

\begin{align*}
  P, Q \;::=\; & \mathbf{0} \\
    \mid\; & P \mid Q \\
    \mid\; & \mathbf{new}\; P \\
    \mid\; & U[\,i_0, \ldots, i_{k-1}\,]\;;\; P \\
    \mid\; & \mathbf{out}\; i \\
    \mid\; & \mathbf{in}\; i\;\{P,\; Q\}
\end{align*}

\paragraph{Shift operation.}
\begin{align*}
  \mathrm{shift}(c,\; \mathbf{new}\; P)
    &= \mathbf{new}\;\mathrm{shift}(c+1,\; P) \\
  \mathrm{shift}(c,\; P \mid Q)
    &= \mathrm{shift}(c,\; P) \mid \mathrm{shift}(c,\; Q) \\
  \mathrm{shift}(c,\; U[\mathit{is}]\;;\; P)
    &= U[\mathrm{shift}(c,\; \mathit{is})]\;;\; \mathrm{shift}(c,\; P) \\
  \mathrm{shift}(c,\; \mathbf{out}\; i)
    &= \mathbf{out}\;\mathrm{shift}(c,\; i) \\
  \mathrm{shift}(c,\; \mathbf{in}\; i\;\{P,\; Q\})
    &= \mathbf{in}\;\mathrm{shift}(c,\; i)\;\{\mathrm{shift}(c,\; P),\;
       \mathrm{shift}(c,\; Q)\}
\end{align*}

\subsection{Denotational Semantics}

These processes act linearly on the distribution:
\[
  \llbracket U[\,i_0, \ldots, i_{k-1}\,]\;;\; P \rrbracket(n, \ket{\psi})
  \;=\;
  \llbracket P \rrbracket(n,\; U\ket{\psi})
\]
where $U\ket{\psi}$ is shorthand for acting on those qubits.
\[
  \llbracket \mathbf{new}\; P \rrbracket(n, \ket{\psi})
  \;=\;
  \llbracket P \rrbracket(n+1,\; \ket{0} \otimes \ket{\psi})
\]

\paragraph{Parallel composition (interleaving semantics).}
\[
  \llbracket \textstyle\prod_i P_i \rrbracket(n, \ket{\psi})
  \;=\;
  \text{random permutation } \sigma\;
  \left(\bigcirc_{\sigma(i)}
    \llbracket P_{\sigma(i)} \rrbracket\right)(n, \ket{\psi})
\]

\[
  \llbracket P \rrbracket(n, \ket{\psi})
  \otimes
  \llbracket Q \rrbracket(m, \ket{\phi})
  \;\longrightarrow\;
  \llbracket P \mid \mathrm{shift}(Q)^n \rrbracket
    (n+m,\; \ket{\psi}\ket{\phi})
\]

\paragraph{Synchronization is measurement.}
\[
  \llbracket \mathbf{out}\; i \mid \mathbf{in}\; i\;\{P,\; Q\} \rrbracket
    (n, \ket{\psi})
  \;=\;
  \Pr(i{=}\ket{0})\;\llbracket P \rrbracket(n, \ket{\psi}_{0/i})
  \;+\;
  \Pr(i{=}\ket{1})\;\llbracket Q \rrbracket(n, \ket{\psi}_{1/i})
\]
where $\{x/i\}$ means we collapse qubit $i$ of $\psi$ to $\ket{x}$.

\subsection{Structural Congruence}

The usual equations hold:
\begin{align*}
  P \mid \mathbf{0} &= P \\
  P \mid Q &= Q \mid P \\
  (P \mid Q) \mid R &= P \mid (Q \mid R) \\
  \mathbf{new}\; P \mid Q &= \mathbf{new}\;(P \mid \mathrm{shift}(1, Q))
\end{align*}

%% ----------------------------------------------------------------
\section{Conclusions and Future Work}
%% ----------------------------------------------------------------

\end{document}
